\section{Introduction}

Critical-infrastructure and public-sector security programs are required to prioritize
remediation by organizational risk rather than raw severity~\cite{nist80053} and to categorize
systems by mission impact~\cite{fips199}. The strongest published prioritization signals do not encode that
mission dimension. The Exploit Prediction Scoring System (EPSS) estimates the probability that
a CVE is exploited in the wild within twelve months of disclosure~\cite{jacobs2021epss}.
Hygiene-augmented scorers add host-level signals such as patch posture and exposure, improving
precision on a hygiene-defined target~\cite{paper4hygieneprio}. Both remain silent on how much
a given host matters to the mission. A queue ordered purely by exploit likelihood and hygiene
can therefore spend a finite remediation budget on poorly-patched but low-value endpoints while
a freshly-patched Criminal Justice Information Services (CJIS) controller waits.

The gap is structural rather than incidental. EPSS scores a CVE; a hygiene score rates a host
as a potential weak link; neither answers the question an authorizing official actually asks,
which is what the organization stands to lose if a particular host is compromised. A
critical-infrastructure or public-sector endpoint fleet is mission-heterogeneous by
construction: a CJIS-bearing domain controller in a
demilitarized zone and a training-room kiosk may carry the same software and the same exploit
exposure, yet a compromise of the first is a reportable incident and a compromise of the second
is an inconvenience. When remediation capacity is the binding constraint, which it almost always
is in practice, ordering the queue by exploit likelihood and hygiene alone leaves mission value
on the table precisely when the budget is too small to absorb it. This is the regime in which a
mission-aware ordering should help most, and it is the regime our study targets.

This paper tests one narrow, falsifiable question: does adding a pre-registered asset context
score to a hygiene-augmented base scorer improve the prioritization of mission-critical
exposure under a fixed remediation capacity, and what does it cost? The study uses a
pre-registered simulation: a synthetic fleet with structural priors taken from public
guidance, real EPSS and Known Exploited Vulnerability (KEV) signal resampled from a frozen
corpus of 203{,}174 CVEs, hypotheses and thresholds locked before any evaluation seed is
inspected, and bias-corrected and accelerated (BCa) bootstrap confidence
intervals~\cite{efron1987bca,efron1994introduction} as the evidentiary standard. The discipline
of fixing hypotheses, metrics, and thresholds in a dated protocol before inspecting the
evaluation data follows the pre-registration practice now standard in empirical software
engineering~\cite{wohlin2012experimentation}, and it is what lets us report the primary result
as partially supported rather than re-cut a threshold to manufacture a clean pass.

\noindent\textbf{Contributions.}
\begin{enumerate}
\item \textbf{CAP-G}, a context-aware scorer that multiplies a hygiene-augmented base score by
a host context factor built from FIPS-199 criticality, network zone, and data sensitivity
(Section~\ref{sec:method}).
\item A \textbf{mission-weighted evaluation} (MWP@K, CWER@K) appropriate to
critical-infrastructure and public-sector fleets
and distinct from the context-blind precision the base scorer optimizes
(Section~\ref{sec:eval}).
\item A \textbf{pre-registered, honestly-reported result} (Section~\ref{sec:results}): a large
gain at scarce capacity that decays as capacity grows, is fully explained by fleet
heterogeneity, costs nothing on the context-blind metric, and misses the capacity-flat success
bar, reported as partial rather than re-cut to a clean pass.
\item An \textbf{exact mechanism control and ablation} (Section~\ref{sec:results}): a
homogeneous-fleet control on which the context factor provably collapses to the base scorer,
isolating fleet heterogeneity as the source of the gain, together with a dimension ablation that
identifies asset criticality as the load-bearing context attribute.
\end{enumerate}

The result is deliberately framed as a bounded, regime-specific finding rather than a universal
improvement. CAP-G helps where a triage budget is scarce and a fleet is mission-heterogeneous,
which is the common operating condition for a public-sector security operations team; it does
nothing where capacity is generous or the fleet is uniform. Stating those bounds plainly, and
holding to a pre-registered bar that the data narrowly miss, is the point of the study rather
than a concession.
